	\begin{center}
		\section{常用 Python 命令}
	\end{center}

\begin{multicols*}{2}

	\subsection{启动交互解释器}

	\textbf{启动命令因环境而异，依次试}：\texttt{python3} / \texttt{python} / \texttt{py}。

	\begin{itemize}
		\item \textbf{Linux 赛场机}：一般只有 \texttt{python3}（裸 \texttt{python} 可能不存在或指向 Python 2）。
		\item \textbf{Windows}：官方安装器装的是 py launcher，敲 \texttt{py}；从 Microsoft Store 装的或勾了 PATH 的则 \texttt{python} 直接可用。
		\item 退出：\texttt{exit()}（或 Unix \texttt{Ctrl-D} / Windows \texttt{Ctrl-Z} + \texttt{Enter}）。
	\end{itemize}

	\subsection{当计算器：科学计数法速估}

	赛场想快速估算（阶乘 / 组合数 / $\log$ / 大整数幂会不会爆 \texttt{ll}）时，进解释器先敲这两行：

	\begin{minted}{python}
>>> from math import *
>>> def p(x): print(f"{x:.2e}")
...
>>> p(factorial(20))
2.43e+18
>>> p(comb(30, 15))
1.55e+08
>>> p(2**64)
1.84e+19
>>> log2(1e9)
29.897352853986263
	\end{minted}

	\texttt{p(x)} 按科学计数法两位小数打印，配合 Python 天然无界的大整数，判「爆不爆 \texttt{ll}（上限 $\approx$ \texttt{9.22e18}）/ 复杂度量级多大」一眼出结果。

\end{multicols*}
